% ===== §5 =====
\section{Structure and Reproduction}
\label{p24:sec:reproduction}

\noindent
The third tradition asks a question the first two do not: not what culture is or how it means, but what culture \emph{does}, and in particular what work it performs in holding a social order together and passing it on. Its two moments, the structural and the reproductive, are distinct, but they share the conviction that beneath the surface of cultural life lies a structure that organizes it, and that the structure is prior to and largely hidden from the agents whose practice it governs.

In its structural moment, associated with L\'evi-Strauss, the thesis is that the diversity of cultural forms---myths, kinship systems, classifications---is the surface transformation of a small number of deep structures, ultimately grounded in the binary operations of the mind, so that the anthropologist's task is to recover the invariant structure of which the visible culture is a permutation \parencite{levistrauss1963}. This is a powerful vision, and it sees something the interpretive tradition can underplay: that culture has a structure with its own logic, not reducible to the meanings its participants avow, a grammar that generates the sayable without being itself said. But it purchases this depth at the price of the agents entirely. The structure operates behind their backs and through them; they are its bearers, not its authors; and the structure itself is timeless, a synchronic order from which history and emergence have been evacuated on principle. Culture here is at its most settled precisely where it is at its most sophisticated: the deep structure is the very type of a thing that is always already there.

Bourdieu's reproductive moment is the one this paper must engage most closely, because it is the sharpest and because its central concept, cultural capital, stands in the most direct opposition to what this paper will claim about intimate culture. Bourdieu set out to overcome the antinomy of structure and agency, and his instrument was the \emph{habitus}: a system of durable, transposable dispositions, laid down in the body by a whole history of conditioning, which generates practices without being either free choice or mechanical determination, an incorporated structure that reproduces the objective structure that produced it \parencite{bourdieu1984}. Around the habitus he built an economy of \emph{capital} in several forms---economic, social, and, decisively, cultural---and showed how taste, comportment, and cultivation, far from being the free play of sensibility they present themselves as, function as \emph{cultural capital}: assets, unequally distributed along class lines, that can be accumulated, that can be converted into economic and social advantage, and that are transmitted across generations through the family and the school so as to reproduce the class structure while concealing that reproduction behind the appearance of individual merit \parencite{bourdieu1986}. This is a genuine and permanent achievement, and it is where the historical-materialist line of this paper first enters the argument. Bourdieu makes visible what the interpretive tradition leaves in shadow: that culture is not an innocent web of meaning but an instrument in the reproduction of social relations, that the power to say what counts as cultivation is a class power, and that the question of who owns and transmits cultural competence is a question of political economy. Any account of culture that forgot this---any celebration of intimate culture as a pure generative idyll---would have earned Bourdieu's contempt, and rightly.

\subsection*{The extractability of cultural capital, and the point of divergence}

But the very concept that makes Bourdieu's analysis so powerful marks the exact place where this paper will part from it, and the point must be stated with precision, because a carelessness here would either concede too much or claim too much. The category \emph{capital} imports into the analysis of culture the defining logic of capital as such: that it is a value which can be accumulated, which can be converted across domains, and which can be made to yield a return. Cultural capital, on Bourdieu's account, is accruable, is convertible into economic and social advantage, and is, in its objectified and institutionalized states above all, extractable---made to work, cashed in, put to the reproduction of advantage. It is true, and Bourdieu is careful about this, that cultural capital in its embodied state cannot be handed over like a coin, since it is incorporated into a habitus over time and dies with its bearer; the point of divergence is not a claim that Bourdieu missed this nuance. The point of divergence is that for Bourdieu even embodied cultural capital, however slowly acquired, is still \emph{capital}: it is oriented to accumulation and return, its value is its convertibility, and its cultural meaning is legible, in the last analysis, as position in a field of struggle. What Bourdieu's frame cannot register---because its founding metaphor forbids it---is a form of culture whose value is precisely \emph{not} of this kind: a culture generated within a relation whose worth is non-accruable, non-convertible, and non-extractable, because it is not an asset that positions its holders in an external field but a real bond that exists only in and for the relation itself. This is the form of culture the intimate dyad generates, and it is the antithesis, term for term, of cultural capital: where cultural capital is transferable in principle and extractable in fact, intimate culture is non-transferable in principle and, this paper will argue, resistant to extraction at its core. The disagreement is recorded here and settled later, at the synthesis and again in the phenomenology of the dyad; for now it is enough to have marked that the strongest theory of culture-as-reproduction is built on a metaphor---capital---whose reach this paper means to bound, by exhibiting a culture that lies wholly outside it.
